
\section*{2. ``Open properties'' of groups and group morphisms
  locally of finite presentation}
\phantomsection\label{sec:sga3-VIB-2}
\addcontentsline{toc}{section}{2. ``Open properties'' of groups and
  group morphisms locally of finite presentation}

\medskip
\phantomsection\label{par:sga3-VIB-2.0}%
\noindent\textbf{2.0.---}
In all that follows, \(S\) will denote an arbitrary scheme; an
\(S\)-group scheme will be called an \(S\)-group. Given an
\(S\)-group \(G\), we shall denote by \(\epsilon\) the identity
section, by
\[
  c:G\longrightarrow G
\]
the inversion morphism, and by
\[
  \mu:G\times_S G\longrightarrow G
\]
the multiplication morphism. For every \(S\)-scheme \(X\), we
shall denote by \(\pi\), or by \(\pi_X\), the structural morphism
\[
  X\longrightarrow S.
\]

Given a property \(\mathscr P(u)\) of a morphism of \(S\)-schemes
\[
  u:X\longrightarrow Y,
\]
we shall say that \(\mathscr P(u)\) is \emph{stable under base
change} if, whenever \(u\) satisfies \(\mathscr P(u)\), the same is
true of the morphism \(u_{(Y')}\), for every \(S\)-morphism
\(Y'\to Y\). We say that \(\mathscr P(u)\) is \emph{local in
nature for the topology \(\mathcal T\)}
(cf.\ Expos\'e~IV, \S\S~4 and 6) if \(\mathscr P(u)\) satisfies the
following two conditions:
\begin{enumerate}
\item[\textup{a)}] \(\mathscr P(u)\) is stable under base change;
\item[\textup{b)}] whenever there is a family of \(S\)-morphisms
  \(Y_i\to Y\), covering for the topology \(\mathcal T\), such
  that each morphism \(u_{(Y_i)}\) satisfies \(\mathscr P(u)\),
  then \(u\) satisfies \(\mathscr P(u)\).
\end{enumerate}

\medskip
\phantomsection\label{prop:sga3-VIB-2.1}%
\noindent\textbf{Proposition 2.1.---}
\emph{Let \(\mathscr P(u)\) be a property of a morphism of
\(S\)-schemes, local in nature for the \(\mathrm{fpqc}\) topology.
Let \(u:G\to H\) be a morphism of \(S\)-groups. Suppose that \(G\)
is flat and universally open over \(S\).}

\emph{Let \(W\) be the largest open subset of \(H\) over which \(u\)
satisfies the property \(\mathscr P(u)\), and let
\(V=u^{-1}(W)\). Then
\[
  U=\pi_G(V)
\]
is open in \(S\), and \(V\) is an open group subscheme of
\(G|_U\).}

The existence of a largest open subset \(W\) of \(H\) over which
\(u\) satisfies \(\mathscr P(u)\) follows from the fact that
\(\mathscr P(u)\) is local in nature for the Zariski topology.
Since \(\pi_G\) is universally open, \(\pi_G(V)\) is an open
subset \(U\) of \(S\). It is enough to show that \(V\) is a group
subscheme of \(G|_U\). We may therefore suppose that \(U=S\).

Set
\[
  G'=G\times_SV,\qquad
  H'=H\times_SV,\qquad
  V'=V\times_SV,\qquad
  W'=W\times_SV,
\]
and \(u'=u_{(V)}\). Let \(W'_1\) be the largest open subset of
\(H'\) over which \(u'\) satisfies \(\mathscr P(u)\). Since \(V\)
is flat and universally open over \(S\), the same is true of
\(H'\) over \(H\), and \hyperref[lem:sga3-VIB-2.1.1]{Lemma~2.1.1} below shows that
\[
  W'_1=W'.
\]
Now consider the automorphism \(a\) (respectively, \(b\)) of the
\(V\)-scheme \(G'\) (respectively, \(H'\)), namely right
translation by the inverse of the diagonal section \(\delta\)
(respectively, by the inverse of \(u(\delta)\)), defined by
\[
  a(g,v)=(g\cdot v^{-1},v),
  \qquad
  \bigl(\text{respectively, }b(h,v)
    =(h\cdot u(v^{-1}),v)\bigr),
\]
for every morphism \(T\to S\), every \(g\in G(T)\),
\(v\in V(T)\), and \(h\in H(T)\).%
\begingroup
\renewcommand{\thefootnote}{\fnsymbol{footnote}}%
\footnote[1]{\label{note:sga3-VIB-2.1-b-formula}Translator's note:
the French source omits \(b\) before \((h,v)\) in the second
displayed definition. The preceding sentence defines the
automorphism \(b\), and the identity immediately below requires
the displayed map to be \(b(h,v)\); the symbol is supplied here.}%
\endgroup{}
It is clear that
\[
  u'\circ a=b\circ u',
\]
which shows that \(W'\) is stable under \(b\), hence that \(V'\)
is stable under \(a\), so that \(V\) is a group subscheme of \(G\).

\medskip
\phantomsection\label{lem:sga3-VIB-2.1.1}%
\noindent\textbf{Lemma 2.1.1.---}
\emph{Let \(\mathscr P(u)\) be a property of an \(S\)-morphism
\(u\), local in nature for the \(\mathrm{fpqc}\) topology.
Consider a cartesian square of morphisms of \(S\)-schemes:}
\begin{center}
  \phantomsection\label{diag:sga3-VIB-2.1.1-cartesian-square}%
  \begin{tikzcd}[row sep=large,column sep=large]
    X' \arrow[r,"f'"] \arrow[d] &
    Y' \arrow[d,"g"'] \\
    X \arrow[r,"f"] &
    Y
  \end{tikzcd}
\end{center}
\emph{where \(g\) is flat and open. Let \(W\) (respectively,
\(W'_1\)) be the largest open subset of \(Y\) (respectively,
\(Y'\)) over which \(f\) (respectively,
\(f'=f_{(Y')}\)) satisfies \(\mathscr P(u)\). Then}
\[
  \emph{\(W'_1=W\times_Y Y'\).}
\]

Set \(W'=W\times_Y Y'\). Since \(\mathscr P(u)\) is stable under
base change, we have
\[
  W'\subset W'_1.
\]
Since \(g\) is open,
\[
  W_1=g(W'_1)
\]
is an open subset of \(Y\). Set
\[
  V_1=f^{-1}(W_1),
  \qquad
  V'_1=V_1\times_{W_1}W'_1;
\]
it is clear that \(V'_1=f'^{-1}(W'_1)\). Since \(g\) is flat and
open, the morphism
\[
  W'_1\longrightarrow W_1=g(W'_1)
\]
induced by \(g\) is faithfully flat and open, hence is a covering
for the \(\mathrm{fpqc}\) topology
(cf.\ \SGARefLink{sga3:iv:proposition:6:3:1:item:iv:01}{Expos\'e~IV, 6.3.1(iv)}). Since the morphism
\[
  V'_1\longrightarrow W'_1
\]
induced by \(f'\) satisfies \(\mathscr P(u)\), the same is true of
the morphism \(V_1\to W_1\) induced by \(f\). Thus
\[
  W_1\subset W,
  \qquad
  W'_1\subset g^{-1}(W_1)
    \subset g^{-1}(W)=W',
\]
and therefore \(W'=W'_1\).

\medskip
\phantomsection\label{rem:sga3-VIB-2.1.2}%
\noindent\textbf{Remark 2.1.2.---}
\emph{A great many properties of a morphism are local in nature
for the \(\mathrm{fpqc}\) topology; among them are the properties
of being flat, (universally) open, (locally) of finite type, of
finite presentation, quasi-finite
(cf.\ EGA~IV\(_2\), 2.5.1, 2.6.1, and 2.7.1), smooth, \'etale,
and unramified (EGA~IV\(_4\), 17.7.3).
}

\emph{The proof of \hyperref[prop:sga3-VIB-2.1]{2.1} in fact uses only base changes by flat morphisms.
The proposition therefore applies to a property satisfying
condition \textup{b)} of \hyperref[par:sga3-VIB-2.0]{2.0} relative to the \(\mathrm{fpqc}\)
topology and stable under base change by flat morphisms (for
example, the property of being quasi-compact and dominant).
}

\emph{Of course, one can state an analogous proposition concerning
properties local in nature for a topology \(\mathcal T\) finer
than the Zariski topology, the condition to be verified by \(G\)
then being that \(\pi_G\) be universally open and covering for
the topology \(\mathcal T\).
}

\emph{In particular, if \(G\) is flat and locally of finite presentation
over \(S\), there is an analogous statement for properties stable
under base change by flat morphisms locally of finite presentation
and satisfying condition \textup{b)} of \SGARefLink{par:sga3-VIB-2.0}{2.0} relative to the
\(\mathrm{fpqc}\) topology (for example, the properties of being
regular, reduced, Cohen--Macaulay, and so on; cf.\
EGA~IV\(_2\), 6.8).}

\medskip
\phantomsection\label{prop:sga3-VIB-2.2}%
\noindent\textbf{Proposition 2.2.---}
\emph{Let \(G\) and \(H\) be two \(S\)-groups, and let
\(u:G\to H\) be a morphism of \(S\)-groups. Then:}%
\footnote{\label{note:sga3-VIB-19}Editor's note: in case
\textup{\SGARefLink{prop:sga3-VIB-2.2:item:i}{(i)}}, the open subset \(V\) consists of all the points of
\(G\) at which \(u\) is smooth, respectively \'etale,
respectively of finite presentation and flat. By contrast, in case \textup{\SGARefLink{prop:sga3-VIB-2.2:item:ii}{(ii)}}, \(V\) is the largest
open subset contained in the set \(E\) of points of \(G\) at
which \(u\) is universally open, but \(E\) is not necessarily
open, as the following example shows
(EGA~IV\(_3\), 14.1.3(i)). Let \(k\) be a field, let
\(H=S=\operatorname{Spec}A\), where \(A=k[x]\), and let \(G\) be
the \(S\)-group \(\operatorname{Spec}A[y]/(xy)\). Then \(E\) is
the identity section of \(G\), which is not open.}
\begin{enumerate}
\item[\textup{\SGAInlineRefTarget{prop:sga3-VIB-2.2:item:i}(i)}] \emph{Suppose that \(G\) or \(H\) is flat over
  \(S\), and that \(G\) or \(H\) is locally of finite presentation
  over \(S\). Let \(V\) be the largest open subset of \(G\) such
  that the restriction of \(u\) to \(V\) is flat and locally of
  finite presentation (respectively, smooth, or \'etale). Then
  \(U=\pi_V(V)\) is open in \(S\), and \(V\) is an open group
  subscheme of \(G|_U\).}
\item[\textup{\SGAInlineRefTarget{prop:sga3-VIB-2.2:item:ii}(ii)}] \emph{Suppose that \(G\) or \(H\) is
  universally open over \(S\), and let \(V\) be the largest open
  subset of \(G\) such that the restriction of \(u\) to \(V\) is
  universally open. Then \(U=\pi_V(V)\) is open in \(S\), and
  \(V\) is an open group subscheme of \(G|_U\).}
\end{enumerate}

We first prove \textup{\SGARefLink{prop:sga3-VIB-2.2:item:i}{(i)}}. We show that the restriction
\(\pi_V\) of \(\pi_G\) to \(V\) is flat and locally of finite
presentation:
\begin{enumerate}
\item[\textup{a)}] if \(\pi_G\) is flat (respectively, locally of
  finite presentation), then so is \(\pi_V\);
\item[\textup{b)}] if \(\pi_H\) is flat (respectively, locally of
  finite presentation), then so is \(\pi_V\), as the composite of
  the restriction of \(u\) to \(V\) with \(\pi_H\).
\end{enumerate}
Thus, in the four cases contemplated in the statement, \(\pi_V\)
is flat and locally of finite presentation, hence universally open
(EGA~IV\(_2\), 2.4.6). Set \(U=\pi_V(V)\); then \(U\) is open in
\(S\). It is therefore enough to show that \(V\) is an open group
subscheme of \(G|_U\), and we may suppose that \(U=S\).

Set
\[
  G'=G\times_SV,\qquad
  H'=H\times_SV,\qquad
  V'=V\times_SV,
  \qquad\text{and}\qquad
  u'=u_{(V)}.
\]
Since \(V\) is flat and locally of finite presentation over \(S\),
the same is true of \(H'\) over \(H\). By EGA~IV\(_4\), 17.7.4,
\(V'\) is then the largest open subset of \(G'\) such that the
restriction of \(u'\) to \(V'\) is flat and locally of finite
presentation (respectively, smooth, or \'etale). With the
automorphisms \(a\) and \(b\) defined as in the proof of \SGARefLink{prop:sga3-VIB-2.1}{2.1}, it
is clear that \(V'\) is stable under \(a\), hence that \(V\) is a
group subscheme of \(G\).

We now prove \textup{\SGARefLink{prop:sga3-VIB-2.2:item:ii}{(ii)}}. The restriction \(\pi_V\) of \(\pi_G\)
to \(V\) is universally open, either because \(\pi_G\) is
universally open, or because it is the composite of the restriction
of \(u\) to \(V\) with \(\pi_H\), when \(\pi_H\) is universally
open. Set \(U=\pi_V(V)\); then \(U\) is open in \(S\). It is enough
to show that \(V\) is an open group subscheme of \(G|_U\), and we
may suppose that \(U=S\).

Set, as before,
\[
  G'=G\times_SV,\qquad
  H'=H\times_SV,\qquad
  V'=V\times_SV,
  \qquad\text{and}\qquad
  u'=u_{(V)}.
\]
Then \(\pi_V:V\to S\) is surjective and universally open; the same
is therefore true of \(G'\to G\). Consequently, \(V'\) is the
largest open subset of \(G'\) such that the restriction of \(u'\)
to \(V'\) is universally open, by EGA~IV\(_3\), 14.3.4(i),(ii).
With the automorphisms \(a\) and \(b\) defined as before, it is
clear that \(V'\) is stable under \(a\), hence that \(V\) is a
group subscheme of \(G\).

\medskip
\phantomsection\label{cor:sga3-VIB-2.3}%
\noindent\textbf{Corollary 2.3.---}
\emph{Let \(G\) be an \(S\)-group, and let \(V\) be the largest
open subset of \(G\) that is flat and locally of finite
presentation (respectively, smooth, \'etale, or universally open)
over \(S\). Then \(U=\pi_V(V)\) is open in \(S\), and \(V\) is an
open group subscheme of \(G|_U\).}

It is enough to apply \hyperref[prop:sga3-VIB-2.2]{2.2} when \(H\) is the identity \(S\)-group
and \(u\) is the unique morphism of \(S\)-groups \(G\to H\), since
\(\pi_H\) is then an isomorphism and
\[
  \pi_G=\pi_H\circ u.
\]

\medskip
\phantomsection\label{cor:sga3-VIB-2.4}%
\noindent\textbf{Corollary 2.4.---}
\emph{Let \(G\) be an \(S\)-group. If there is a neighborhood \(X\)
of the identity section having one (or more) of the following
properties:}
\begin{center}
  \emph{\(X\) is flat and locally of finite presentation
  (respectively, smooth, \'etale, or universally open) over \(S\),}
\end{center}
\emph{then there is an open group subscheme of \(G\) having the
same properties.}

It is enough to apply \hyperref[cor:sga3-VIB-2.3]{2.3}, observing that here, with the notation
of \hyperref[prop:sga3-VIB-2.2]{2.2}, we have
\[
  \epsilon(S)\subset V,
\]
and hence \(U=S\).
